\section{The exterior algebra model in coordinates}\label{app:exterior}

We record the computation behind \Cref{thm:fmtheta} in coordinates, both
because it is the only place in the paper where a sign has to be pinned down
and because the model is the one the verification script implements.

Let $g\ge1$ and let $W$ be the $\QQ$-vector space with basis
$e_{1},\dots,e_{2g},f_{1},\dots,f_{2g}$. We identify
$V=\langle e_{1},\dots,e_{2g}\rangle$ with $H^{1}(A,\QQ)$ and
$V^{\vee}=\langle f_{1},\dots,f_{2g}\rangle$ with $H^{1}(\Av,\QQ)$, so that
$\bigwedge^{*}W=H^{*}(A\times\Av,\QQ)$.

In $\bigwedge^{*}W$ a basis is indexed by subsets $S$ of
$\{1,\dots,4g\}$, with the convention that the generator numbered $i$ is
$e_{i}$ for $i\le2g$ and $f_{i-2g}$ for $i>2g$, and that the basis element
$w_{S}$ is the wedge of the generators of $S$ in increasing order. The product
is $w_{S}w_{T}=\varepsilon(S,T)\,w_{S\cup T}$ when $S\cap T=\emptyset$ and zero
otherwise, where
\begin{equation}\label{eq:shufflesign}
  \varepsilon(S,T)=(-1)^{\#\{(s,t)\in S\times T\,:\,s>t\}} .
\end{equation}

The three distinguished elements are
\[
  \ell=\sum_{i=1}^{2g}e_{i}\wedge f_{i},
  \qquad
  \theta=\sum_{j=1}^{g}e_{j}\wedge e_{g+j},
  \qquad
  \theta'=\sum_{j=1}^{g}f_{j}\wedge f_{g+j},
\]
and $p_{*}$ is the linear map $\bigwedge^{*}W\to\bigwedge^{*}V$ that sends
$w_{S}$ to $\varepsilon(S\cap[1,2g],\,S\setminus[1,2g])\,w_{S\cap[1,2g]}$ when
$S\supseteq\{2g+1,\dots,4g\}$ and to zero otherwise.

\begin{lemma}\label{lem:appsign}
Let $0\le k\le g$, let $T\subseteq\{1,\dots,g\}$ with $|T|=k$, let
$U=\{1,\dots,g\}\setminus T$ and $S=U\cup(g+U)$. In the product
$\theta'^{\,k}\cdot e^{\ell}$ the unique term contributing $w_{S}$ after
$p_{*}$ is
\[
  k!\;\Bigl(\prod_{j\in T}f_{j}f_{g+j}\Bigr)\cdot\Bigl(\prod_{i\in S}e_{i}f_{i}\Bigr),
\]
and its image under $p_{*}$ is
$k!\,\sigma(g,k)\,\prod_{j\in U}e_{j}\wedge e_{g+j}$ with
$\sigma(g,k)=(-1)^{g(g+1)/2+k}$, independent of $T$.
\end{lemma}

\begin{proof}
Uniqueness of the term and the identification of $S$ are established in the
proof of \Cref{thm:fmtheta}. The sign is the product of the two shuffle signs
\eqref{eq:shufflesign} relating the displayed word to the target word
$\bigl(\prod_{j\in U}e_{j}e_{g+j}\bigr)\bigl(f_{1}\cdots f_{2g}\bigr)$.
Relabelling the symplectic pairs by a permutation $\pi$ of $\{1,\dots,g\}$
carries the word for $T$ to the word for $\pi(T)$ and multiplies both the
displayed and the target word by the sign of the induced permutation of
generators, so the quotient is unchanged; hence $\sigma$ depends only on
$(g,k)$. Evaluating at $T=\{1,\dots,k\}$ gives the stated value, and the
evaluation is carried out for all $g\le12$ and all $k$ in
\Cref{app:scripts}, item~(II).
\end{proof}

\begin{remark}
The computation is small enough to do by hand at $g=2$. There
$\ell=e_{1}f_{1}+e_{2}f_{2}+e_{3}f_{3}+e_{4}f_{4}$,
$\theta=e_{1}e_{3}+e_{2}e_{4}$ and $\theta'=f_{1}f_{3}+f_{2}f_{4}$, and one
finds $\Phi_{\cP}(1)=-\tfrac12\theta^{2}$, $\Phi_{\cP}(\theta')=\theta$ and
$\Phi_{\cP}(\theta'^{2})=-2$, in agreement with \eqref{eq:fmtheta} for $g=2$,
where $(-1)^{g(g+1)/2}=(-1)^{3}=-1$.
\end{remark}
